% Measured variant of the page-1 framing paragraph.
% For task t-mozveh1b (Mathematics) and t-mozv8oul (Biological Science).
% Same content, lower-temperature prose — pairs with the aggressive variant so
% editors can pick one. References Math Theorem 2.1 and the exchangeability bound (Lemma S1).

\paragraph*{The limits of label-based evaluation.}
Standard evaluation of single-cell integration relies on quantities---$\mathrm{ARI}$, $\mathrm{NMI}$, $\mathrm{ASW}$, and related scores in scIB, scIB-E, and RBET---that are functionals of annotator labels. A formal consequence, stated as Theorem~\ref{thm:labelblind}, is that every label-based functional is invariant under a class of local coupling deformations that can eliminate an unannotated rare component by folding it onto an adjacent annotated population. In other words, the failure mode the field has most worried about for at least five years---loss of rare subpopulations that no atlas has yet named---sits outside the measurable sets of its own evaluation apparatus. The failure is not that existing metrics are imprecise; it is that they cannot, in principle, observe the events of interest. This matters downstream: the inference ``integrated data are safe for discovery because no label-based metric dropped'' conflates \emph{absence of evidence} with \emph{evidence of absence} along a dimension the evaluation framework cannot read. A productive response is not to reject integration but to evaluate it with instruments that do not presuppose the annotation. Accordingly, we propose REAL (Rare-subpopulation Erasure under Alignment, Label-free), which assesses integration through a label-free, geometrically-defined notion of local-mass preservation and is calibrated by a held-out known-rare protocol under an explicit exchangeability bound (Lemma~\ref{lem:exchangeability}). REAL does not settle questions of biological novelty; it marks where the evaluation framework has been blind.
